\documentclass[11pt,a4paper]{article}

% --- Pakete ---
\usepackage[ngerman]{babel}
\usepackage[utf8]{inputenc}
\usepackage[T1]{fontenc}
\usepackage{lmodern} % Computer Modern (empfohlen)
\usepackage{geometry}
\geometry{
  a4paper,
  left=2.5cm,
  right=2.5cm,
  top=2.5cm,
  bottom=2.5cm
}
\usepackage{setspace}
\setstretch{1.24} % entspricht Empfehlung (1–1.5 Zeilenabstand)
\usepackage{microtype}
\usepackage{csquotes}
\usepackage{amsmath,amssymb}
\usepackage{graphicx}
\usepackage{caption}
\usepackage{subcaption}
\usepackage{siunitx}     % saubere Einheiten, Einheiten aufrecht
\sisetup{detect-all}
\usepackage{booktabs}
\usepackage{hyperref}
\usepackage{fancyhdr}

% --- Kopf- und Fußzeile ---
\pagestyle{fancy}
\fancyhf{}
\fancyhead[L]{Abschnitt \leftmark}
\fancyfoot[C]{\thepage}

% --- Nummerierte Gleichungen ---
\numberwithin{equation}{section}

% --- Abbildungs- und Tabellenbeschriftungen ---
\captionsetup{font=small,labelfont=bf}

% Kopfzeile: links aktueller Abschnitt
\fancyfoot[C]{\thepage}

% Section-Mark definieren, damit \leftmark das Richtige enthält
\renewcommand{\sectionmark}[1]{%
  \markboth{\thesection\ #1}{}%
}

% --- Metadaten (für Titelseite) ---
\newcommand{\Lehrveranstaltung}{Praktikum Physikalische Eigenschaften der Materialien}
\newcommand{\Universitaet}{Universität Augsburg}
\newcommand{\Semester}{Wintersemester 2025/2026}
\newcommand{\Versuchsnummer}{2}
\newcommand{\Versuchstitel}{Solarzelle}
\newcommand{\Gruppe}{C}
\newcommand{\Versuchsdatum}{02.03.2026}
\newcommand{\Abgabedatum}{TT.MM.JJJJ}
% Bei Einzelabgabe:
%\newcommand{\Abgabeart}{Versuchsprotokoll}
%\newcommand{\Autor}{Max Mustermann}
% Bei gemeinsamer Abgabe:
\newcommand{\Abgabeart}{Gemeinsames Versuchsprotokoll}
\newcommand{\Autor}{Fabio Piskorz, Fabio Nogielsky, Konstantin Szantho von Radnoth }

% ---------------------------------------------------
% Titelseite (ohne Seitenzahl)
% ---------------------------------------------------
\begin{document}
\pagenumbering{gobble} % keine Seitenzahl auf Titelseite

\begin{titlepage}
    \centering
    {\Large \Lehrveranstaltung\par}
    \vspace{0.5cm}
    {\large \Universitaet\par}
    \vspace{0.5cm}
    {\large \Semester\par}
    \vspace{1.5cm}
    {\LARGE \Abgabeart\par}
    \vspace{0.5cm}
    {\Large \Versuchsnummer: \Versuchsnummer\par}
    \vspace{0.5cm}
    {\Large \Versuchstitel: \Versuchstitel\par}
    \vspace{1.5cm}
    {\large Gruppe \Gruppe\par}
    \vspace{0.5cm}
    {\large Autor(en): \Autor\par}
    \vspace{1.5cm}
    {\large Versuchsdatum: \Versuchsdatum\par}
    \vspace{0.5cm}
    {\large Abgabedatum: \Abgabedatum\par}

    \vfill
\end{titlepage}


%Inhaltsverzeichnis
\tableofcontents
\newpage
% ab hier Seitenzahlen normal
\pagenumbering{arabic}

% ---------------------------------------------------
% 1 Einleitung (max. ca. 1/2 Seite)
% ---------------------------------------------------
\section{Einleitung}
% Kurz und prägnant, max. 2/3 Seite. [file:1]
% - Warum ist der Versuch interessant/wichtig?
% - Welcher theoretische Hintergrund?
% - Was soll konkret untersucht werden?
Die globale Energieversorgung ist f\"ur einen Gro{\ss}teil der Treibhausgasemissionen verantwortlich. Rund 80 \% des globalen Energiebedarfs werden aus fossilen Energietr\"agern gewonnen \cite{statis1}. Um der Klimakrise entgegenzuwirken und eine langfristig nachhaltige Energieversorgung zu gew\"ahrleisten werden Erneuerbare Energiequellen immer wichtiger. Im Vergleich der vier leistungsst\"arksten Quellen erneuerbarer Energie (Solar-, Wasser-, Windenergie und Biomasse) zeichnet sich ein deutlich Trend fort: Der Ausbau der installierten Leistung nimmt \"uber alle hinweg zu, den gr\"o{\ss}te Anstieg verzeichnet die Solarenergie \cite{iea2024}. \newline
Solarzellen k\"onnen die Strahlungsleistung der Sonne direkt nutzbar machen. Abhängig von Standort und Atmosphäre liefert die Solarstrahlung in Deutschland 900 bis 1200 Kilowattstunden pro Quadratmeter und Jahr \cite{weltphysik}. \newline
In diesem Versuch wird die Erzeugung elektrischer Energie mittels Solarzellen untersucht. Dazu geh\"oren die Bestimmung des Wirkungsgrades, das Auftragen der irreversiblen Verluste in einer U-I-Kennlinie, sowie die Untersuchung der Empfindlichkeit f\"ur unterschiedliche Wellenl\"angen.
Abschlie{\ss}end soll in einem kleinen Beispiel gezeigt werden, wie die gewonnene Energie genutzt werden kann.
\newpage

% ---------------------------------------------------
% 2 Theoretische Grundlagen (1–2 Seiten)
% ---------------------------------------------------
\section{Theoretische Grundlagen}
% Nur das, was für Verständnis und Auswertung nötig ist. [file:1]
% Herleitungen ggf. gekürzt und mit Literaturangaben.

In diesem Abschnitt werden die für den Versuch relevanten physikalischen Größen und Zusammenhänge erläutert. [file:1]

Als Beispiel für eine Gleichung mit erklärten Variablen:
\begin{equation}
    \Delta T = T_{\text{Raum}} - T_{\text{H2O}}.
\end{equation}
Hier bezeichnet \(T_{\text{Raum}}\) die Raumtemperatur und \(T_{\text{H2O}}\) die Wassertemperatur. [file:1]
\newpage

% ---------------------------------------------------
% 3 Versuchsbeschreibung
% ---------------------------------------------------
\section{Versuchsbeschreibung}
% Versuchsaufbau und -durchführung so beschreiben,
% dass der Versuch ohne Anleitung nachvollziehbar ist. [file:1]

In diesem Abschnitt werden der Versuchsaufbau und die Versuchsdurchführung detailliert beschrieben, sodass der Versuch ohne zusätzliche Anleitung reproduzierbar ist. [file:1]
% Beispiel für eine Abbildung mit eindeutiger Nummer:
\begin{figure}[h]
    \centering
    %\includegraphics[width=0.7\textwidth]{aufbau.png}
    \caption{Schematische Darstellung des Versuchsaufbaus.}
    \label{fig:aufbau}
\end{figure}

Im Text wird auf die Abbildung \ref{fig:aufbau} verwiesen und deren Inhalt beschrieben. [file:1]
\newpage
% ---------------------------------------------------
% 4 Auswertung und Diskussion
% ---------------------------------------------------
\section{Auswertung und Diskussion}
% Auswertung nach Aufgabenstellung, inkl. Fehlerrechnung
% und ggf. Vergleich mit Literaturwerten. [file:1]

In diesem Abschnitt werden die Messdaten ausgewertet, Unsicherheiten bestimmt und die Ergebnisse mit Literaturwerten verglichen. [file:1]

% Beispiel-Tabelle:
\begin{table}[h]
    \centering
    \caption{Beispielhafte Messwerte und berechnete Größen.}
    \label{tab:messwerte}
    \begin{tabular}{lccc}
        \toprule
        Messung & Zeit / \si{\second} & Weg / \si{\meter} & Geschwindigkeit / \si{\meter\per\second} \\
        \midrule
        1 & 1.0 & 0.50 & 0.50 \\
        2 & 2.0 & 1.00 & 0.50 \\
        \bottomrule
    \end{tabular}
\end{table}

Die Tabelle \ref{tab:messwerte} zeigt beispielhaft Messwerte mit zugehörigen Einheiten. [file:1]
\newpage
% ---------------------------------------------------
% 5 Zusammenfassung (kurz, max. 1/2 Seite)
% ---------------------------------------------------
\section{Zusammenfassung}
% Kurzfassung des Versuchs und der wichtigsten Ergebnisse,
% typischerweise <= 1/2 Seite. [file:1]

Hier werden die wichtigsten Ergebnisse und Erkenntnisse des Versuchs kurz zusammengefasst und ggf. ein Ausblick gegeben. [file:1]
\newpage
% ---------------------------------------------------
% 6 Laborbuch (Anhang)
% ---------------------------------------------------
\appendix
\newpage
\section{Anhang}
% Digitale Kopie des Laborbuch-Eintrags (z.B. Scan oder Foto). [file:1]

Fügen Sie hier eine gut lesbare digitale Kopie der Laborbuchseiten ein, auf denen die Messdaten dokumentiert sind. [file:1]
\newpage
% ---------------------------------------------------
% 7 Literaturverzeichnis
% ---------------------------------------------------
\section{Literaturverzeichnis}
% Einträge in der Reihenfolge des ersten Auftretens im Text. [file:1]
% Beispiel im IEEE-Stil:

\begin{thebibliography}{9}

\bibitem{Demtroeder1}
W.~Demtröder, \emph{Experimentalphysik 1: Mechanik und Wärme}, Springer, 2003. [file:1]

\bibitem{statis1}
Bundesamt, ``Erneuerbare Energien global'', Destatis, \url{https://www.destatis.de/DE/Themen/Laender-Regionen/Internationales/Thema/umwelt-energie/energie/ErneuerbareEnergienGlobal.html}. [Online]. Abgerufen:  März 2026.

\bibitem{iea2024}
International Energy Agency (IEA), \emph{World Energy Outlook 2024}, Paris, Frankreich: IEA, 2024. [Online]. Verfügbar: \url{https://iea.blob.core.windows.net/assets/efba04ce-2472-4df1-9873-fb170c055259/WorldEnergyOutlook2024.pdf}. [Abgerufen: März 2026.]

\bibitem{weltphysik}
``Sonnenenergie'', Welt der Physik, \url{https://www.weltderphysik.de/gebiet/technik/energie/solarenergie/sonnenenergie/}. [Online]. Abgerufen: März 2026.

% Weitere Einträge hier...


% weitere Quellen nach Bedarf

\end{thebibliography}

\end{document}